\documentclass[10pt]{beamer}

% ===== Essential Packages for pdfLaTeX =====
\usepackage[utf8]{inputenc} % Handles standard text encoding
\usepackage[T1]{fontenc}    % Handles standard font encoding (replaces fontspec)

% 1. Set the text font
\usepackage{lmodern} 

% 2. Load math support
\usepackage{amsmath}
\usefonttheme[onlymath]{serif}

% 3. Load the APS-style math font LAST to override defaults
\usepackage[varg]{newtxmath} 

% 4. Load physics and other utilities
\usepackage{physics}
\usepackage{graphicx}
\usepackage{tikz}
\usepackage{quantikz}
\usetikzlibrary{positioning, shapes.geometric, arrows.meta, shapes.multipart}
\usepackage{booktabs}
\usepackage{listings} 
\usepackage{caption}
\usepackage{hyperref}
\usepackage{xcolor}

\definecolor{periwinkle}{RGB}{204, 204, 255}
\definecolor{periwinkleDark}{RGB}{102, 102, 204}

% ===== Progress Bar Implementation =====
\makeatletter
\setbeamertemplate{footline}{%
    \begin{tikzpicture}[remember picture,overlay]
        \fill[gray!30] (current page.north west) -- ++(\paperwidth,0) -- ++(0,-0.1cm) -- ++(-\paperwidth,0) -- cycle;
        \fill[periwinkleDark] (current page.north west) -- ++(\paperwidth*\insertframenumber/\inserttotalframenumber,0) -- ++(0,-0.1cm) -- ++(-\paperwidth*\insertframenumber/\inserttotalframenumber,0) -- cycle;
        \node[anchor=north east, xshift=-0.2cm, yshift=-0.15cm] at (current page.north east) {\scriptsize \insertframenumber/\inserttotalframenumber};
    \end{tikzpicture}%
}
\makeatother

% Custom Equation Command
\newcommand{\eq}[1]{%
  \begin{equation*}
    \begin{aligned}
      #1
    \end{aligned}
  \end{equation*}
}

\renewcommand{\figurename}{Fig}
\setbeamertemplate{caption}[numbered]

% ===== Theme Settings =====
\usetheme{default}
\usecolortheme{seahorse}
\setbeamercolor{title}{bg=periwinkle}
\setbeamercolor{frametitle}{bg=periwinkle}
\setbeamertemplate{navigation symbols}{}

% ===== Title Page Data =====
\title{Presentation Title}
\subtitle{Presentation Subtitle and Description}
\author{John Doe}
\institute{Affiliation/Company}
\date{\today}

\begin{document}

% --- 1. Title Page ---
\begin{frame}[plain]
    \titlepage
    \centering
    \includegraphics[height=1.5cm]{example-image-a}
\end{frame}

\begin{frame}{Outline}
    \begin{enumerate}
        \item Introduction
        \begin{enumerate}\color{blue!80!black}
            \item Topic \#1
            \item Topic \#2
        \end{enumerate}
        \item Methodology
        \begin{enumerate}\color{blue!80!black}
            \item Topic \#1
            \item Topic \#2
        \end{enumerate}
        \item Conclusion and Outlook
        \begin{enumerate}\color{blue!80!black}
            \item Topic \#1
            \item Topic \#2
        \end{enumerate}
    \end{enumerate}
\end{frame}


% --- 2. All Text Slide ---
\begin{frame}{Introduction to Non-Markovian Noise}
    \textbf{Core Concepts:}
    \begin{itemize}
        \item Quantum systems interact with their environment, leading to decoherence.
        \item Markovian noise assumes the environment has no memory (information flows strictly from system to environment).
        \item \textbf{Non-Markovian noise} involves information backflow, where the environment retains memory of the system's past states.
    \end{itemize}
    \vspace{0.5cm}
    \textbf{Implications for Error Correction:}
    \begin{enumerate}
        \item Standard Quantum Error Correction (QEC) assumes independent, memoryless errors.
        \item Memory effects can cause correlated errors over time, degrading threshold theorems.
    \end{enumerate}
\end{frame}

% --- 3. Text (LHS) and Figure (RHS) ---
\begin{frame}{Experimental Setup}
    \begin{columns}[T] % [T] aligns columns at the top
        
        \begin{column}{0.5\textwidth}
            \textbf{System Overview:}
            \begin{itemize}
                \item The setup involves a superconducting qubit coupled to a readout resonator.
                \item We induce memory effects by strongly coupling the qubit to a secondary lossy cavity.
                \item Measurements are taken using standard Ramsey interferometry.
            \end{itemize}
        \end{column}
        
        \begin{column}{0.5\textwidth}
            \centering
            \includegraphics[width=\textwidth]{example-image-b}
            \captionof{figure}{\scriptsize Schematic of the qubit-cavity architecture.}
        \end{column}
        
    \end{columns}
\end{frame}

% --- 4. Only Figure Slide ---
\begin{frame}{State Tomography Results}
    \centering
    \textbf{Fidelity Decay Over Time} \\
    \vspace{0.3cm}
    \includegraphics[width=0.85\textwidth, height=0.55\textheight]{example-image-c}
    \vspace{0.2cm}
    \captionof{figure}{\small The revivals in fidelity clearly indicate information backflow characteristic of non-Markovian dynamics.}
\end{frame}

% --- 5. Text and Equation Slide ---
\begin{frame}{Mathematical Modeling of the Bath}
    To describe the evolution of the open quantum system, we use the time-local master equation rather than the standard Lindblad form.
    
    \vspace{0.3cm}
    \textbf{Time-Convolutionless (TCL) Master Equation:}
    \eq{
        \frac{d}{dt}\rho(t) &= \mathcal{K}(t)\rho(t) \\
        \mathcal{K}(t)\rho(t) &= -i [H(t), \rho(t)] + \sum_i \gamma_i(t) \left( L_i \rho(t) L_i^\dagger - \frac{1}{2}\{L_i^\dagger L_i, \rho(t)\} \right)
    }
    
    \vspace{0.3cm}
    \begin{itemize}
        \item Here, $\gamma_i(t)$ represents the time-dependent decay rates.
        \item A signature of non-Markovianity is when $\gamma_i(t) < 0$ for some interval of time, indicating information flowing back into the system.
    \end{itemize}
\end{frame}

% --- 6. Text, Equation, and Figure ---
\begin{frame}{Information Measure}
    \begin{columns}[T]
        
        \begin{column}{0.5\textwidth}
            \textbf{Trace Distance:} \\
            A common measure for distinguishability between two quantum states $\rho_1$ and $\rho_2$ is the trace distance:
            \eq{
                &D(\rho_1, \rho_2) = \frac{1}{2} \norm{\rho_1 - \rho_2}_1 \\
                &= \frac{1}{2} \Tr \left( \sqrt{(\rho_1 - \rho_2)^\dagger (\rho_1 - \rho_2)} \right)
            }
            \vspace{0.1cm}
            \footnotesize
            When $\frac{d}{dt}D > 0$, the states become more distinguishable, which is our signature for non-Markovianity.
        \end{column}
        
        \begin{column}{0.5\textwidth}
            \centering
            \vspace{1cm} % Push image down to align nicely with text
            \includegraphics[width=0.8\textwidth]{example-image-a}
            \captionof{figure}{\scriptsize Trace distance oscillation over time.}
        \end{column}
        
    \end{columns}
\end{frame}

% --- 7. Citations / References ---
\begin{frame}{References}
    \vspace{0.2cm}
    \begin{enumerate}
        \item[\textbf{[1]}] Breuer, H. P., Laine, E. M., \& Piilo, J. (2009). \textit{Measure for the degree of non-Markovian behavior of quantum processes in open systems}. Physical Review Letters, 103(21), 210401.
        \vspace{0.2cm}
        \item[\textbf{[2]}] Reference Two (e.g., specific QEC paper).
    \end{enumerate}
\end{frame}

% --- 8. Thank You / Closing ---
\begin{frame}[plain]
    \centering
    \vspace{2cm}
    {\Huge \textbf{Thank You!}}
    
    \vspace{1cm}
    {\Large Questions?}
    
    \vfill
    \scriptsize \textcolor{gray}{Contact: john.doe@mailing.ca}
\end{frame}

\end{document}
